% APPENDIX. Does not count toward the 8-page content limit (verified from the ARR CFP).
% Everything here is reproducible from the artifacts named in each section.
% TODO still to port from paper_aaai2027/main.tex supplementary:
%   - Unfiltered evaluation on n=250
%   - Per-disability-category stratification
%   - Mechanistic accounts of baselines (propagation metrics)
%   - Phase 1 vocabulary threshold robustness
%   - Why SVD, not PCA or ICA
%   - Qualitative examples (6 prompt types) and the cross-method comparison table

\section{Judge protocol}
\label{app:judge}

Responses are scored by Llama-3.1-8B-Instruct against a ten-point accessibility rubric,
with \texttt{temperature}$=0.1$, seed 42, and responses truncated to 2{,}000 characters.
The judge is asked for a line of the form \texttt{Overall Score: X/10}, and a response
whose output cannot be parsed, or which is empty, receives $5.0$.

That default deserves scrutiny, since $5.0$ is also a plausible middling score. Across
the 2{,}500 unsteered responses, none receives exactly $5.0$. In this run the value marks
failure rather than mediocrity. Recoding every $5.0$ to zero moves SADI from $3.19$ to
$1.37$, because $36.4\%$ of its responses are unparseable, and moves no other method by
more than $0.05$. The default therefore understates the collapse, and every quality drop
reported in the main text is conservative.

Scoring is deduplicated by response hash: 21{,}750 responses reduce to 19{,}888 unique
texts, each scored once.

\section{Statistical procedures}
\label{app:stats}

The unit of analysis is one sweep point, meaning one method at one configuration, with
$n{=}250$ responses. Intervals on a mean judge score are percentile bootstrap intervals
over 2{,}000 resamples of the 250 responses, seed 42.

Judge scores take eight distinct values and are strongly left-skewed at degenerate
points, so normality fails and we use the Mann-Whitney $U$ test with rank-biserial
correlation as the effect size rather than a $t$-test. Comparisons of the five published
operating points against the pooled unsteered baseline are Holm-corrected across those
five tests. The Pareto front is a descriptive geometric construction and carries no
$p$-value.

Responses within a sweep point come from different prompts and are treated as
independent. The same prompt recurs across sweep points, so comparisons between points
are not strictly independent; we report them as unpaired regardless, which widens the
intervals rather than narrowing them.

Two consequences of the measurement are worth stating. The judge emits eight distinct
values, so it does not grade fine gradations within a single response; at $n{=}250$ the
resulting interval on a point mean has a median half-width of $0.18$. And our CAD runner
records medicalization in aggregate rather than per response, so the CAD points in
Figure~\ref{fig:frontier} carry no interval on the horizontal axis.

The unpaired and matched designs agree on the comparison they share. The frontier gives
$+0.30$ with interval $[0.07, 0.53]$ for the projection against CAA; the matched-item
test in Section~\ref{sec:paired} gives $+0.35$ with $[0.18, 0.53]$.

\section{Sensitivity of the degeneracy criterion}
\label{app:thresholds}

The criterion in Section~\ref{sec:frontier} uses three cut-offs: uniqueness below
$0.50$, parse-failure rate above $0.10$, and medicalization below $-0.05$. These are
conventional rather than derived, so we varied all three. Across the 27 combinations of
$\{0.3, 0.5, 0.7\} \times \{0.05, 0.10, 0.20\} \times \{-0.10, -0.05, 0\}$ the Pareto
front holds the same twelve configurations, ten of them CAD operators, and the same four
configurations reaching at least half the baseline reduction.

Relative to applying no criterion at all, flagging removes exactly one configuration
from the front, FairSteer at $\alpha{=}8$, which reaches medicalization $0.000$ with 68
of its 250 responses distinct.

\section{Contrast-set controls in full}
\label{app:controls}

Section~\ref{sec:geometry} reports the homogeneity ladder in the main text and summarises
three further controls. This section gives those three, all on the AccessEval contrastive
pool of 2{,}083 pairs on Llama-3.1-8B-Instruct.

\paragraph{Cluster versus pair sampling.} AccessEval expands 234 questions, each listed twice, into
4{,}164 pairs by substituting nine disability categories into templated queries, so pairs are
nested within base queries. Sampling 13 pairs uniformly gives $\mathrm{EVR}_1{=}21.5\%$;
sampling 13 pairs by drawing whole base queries gives $72.2\%$. Same data, same $n$, same
layer, a $3.4$-fold swing from the sampling structure alone. The gap closes as $n$ grows:
$15.6\%$ against $40.4\%$ at $n{\approx}30$, $11.1\%$ against $14.1\%$ at $n{\approx}250$,
and $10.4\%$ against $10.8\%$ at $n{\approx}1000$.

\paragraph{A permuted-pairing null.} We match each clean activation to a different
corrupted one, destroying the contrast while preserving both marginals. At $n{=}13$ the
real data give $\mathrm{EVR}_1{=}21.5\%$ and the null gives $26.2\%$, so the null is
higher. The two curves separate only above $n \approx 250$
(Figure~\ref{fig:controls}). Any concentration estimated from a dozen pairs is consistent
with no contrastive structure at all, which is the regime the reported Target A profile
sits in.

\paragraph{Cluster bootstrap.} Resampling the 2{,}083 pairs of the filtered pool
directly gives $k@80\%{=}40$ with a $95\%$ interval of $[40, 41]$. The pairs are not
independent: every pair is one of several disability variants of a question, and AccessEval
lists each question twice, so the pool comes from only 146 distinct questions. Resampling
those 146 questions instead gives $k@80\%{=}30.1$ $[27, 33]$ at L21, $28.2$ $[26, 30]$ at
L14, and $10.2$ $[9, 11]$ at L0 (150 resamples each). The pair-level and question-level
intervals do not overlap at L21, so the pair-level procedure is biased as well as
over-precise.

\begin{figure*}[t]
\centering
\includegraphics[width=\textwidth]{figure-03-contrast-set-controls.pdf}
\caption{Spectral concentration (left) and prescribed rank (right) against the number of
contrastive pairs, on AccessEval. Solid lines sample pairs uniformly; dashed lines sample
whole base queries; the dotted grey line is a null built by permuting the pairing, which
destroys the contrast while preserving both marginals. Bands are $95\%$ intervals over 40
draws. Below $n \approx 250$ the real data and the null are not separated, so spectra
estimated from a dozen pairs carry no contrastive signal. The dash-dotted line marks the
profile the paper reports for Target A.}
\label{fig:controls}
\end{figure*}


\section{Reproduction}
\label{app:repro}

Generation and judging ran on a single A100 80GB. The five sweeps took approximately
eleven GPU-hours in total and the judging pass approximately three. The diagnostic
itself is cheap next to either: collecting activations over the contrastive pairs is one
forward pass, and the truncated SVD of the resulting $n \times d$ difference matrix takes
seconds. The projection is then applied at $\alpha{=}1$ with no strength search, which is
the practical difference from the additive baselines, each of which requires a
layer-by-strength grid where every cell is a full generation-and-scoring pass.

\section{Complete sweep}
\label{app:sweep}

Tables~\ref{tab:allsweep1}--\ref{tab:allsweep3} list every sweep point behind
Figure~\ref{fig:frontier}: 87 configurations, $n{=}250$ responses each, 21{,}750
responses in total. Quality is the mean judge score with a percentile bootstrap $95\%$
interval over 2{,}000 resamples. ``Uniq'' is the fraction of the 250 responses that are
distinct and ``Fail'' the fraction the judge could not parse. A point is marked
degenerate ($\bullet$) when uniqueness falls below $0.50$, the failure rate exceeds
$0.10$, or medicalization falls below $-0.05$.

\begin{table*}[t]
\centering\small
\begin{tabular}{llrrrrc}
\toprule
Method & Configuration & Med. & Quality \scriptsize[95\% CI] & Uniq & Fail & Deg. \\
\midrule
Angular & max\_sim\_L23/mode\_0/angle\_150 & $-0.222$ & $6.06$ \scriptsize$[5.86,\,6.24]$ & $1.00$ & $0.01$ & $\bullet$ \\
Angular & max\_sim\_L23/mode\_1/angle\_150 & $-0.210$ & $6.04$ \scriptsize$[5.83,\,6.23]$ & $1.00$ & $0.01$ & $\bullet$ \\
Angular & max\_sim\_L23/mode\_0/angle\_120 & $-0.169$ & $6.19$ \scriptsize$[6.00,\,6.38]$ & $1.00$ & $0.00$ & $\bullet$ \\
Angular & max\_sim\_L23/mode\_1/angle\_120 & $-0.154$ & $6.39$ \scriptsize$[6.22,\,6.58]$ & $1.00$ & $0.00$ & $\bullet$ \\
Angular & max\_sim\_L23/mode\_0/angle\_180 & $-0.071$ & $6.02$ \scriptsize$[5.82,\,6.21]$ & $1.00$ & $0.02$ & $\bullet$ \\
Angular & max\_sim\_L23/mode\_1/angle\_180 & $-0.070$ & $6.34$ \scriptsize$[6.15,\,6.52]$ & $1.00$ & $0.01$ & $\bullet$ \\
Angular & max\_norm\_L31/mode\_0/angle\_150 & $-0.060$ & $6.17$ \scriptsize$[5.98,\,6.35]$ & $1.00$ & $0.00$ & $\bullet$ \\
Angular & max\_norm\_L31/mode\_0/angle\_180 & $-0.052$ & $5.53$ \scriptsize$[5.32,\,5.74]$ & $1.00$ & $0.00$ & $\bullet$ \\
Angular & max\_norm\_L31/mode\_1/angle\_180 & $-0.052$ & $5.66$ \scriptsize$[5.44,\,5.86]$ & $1.00$ & $0.00$ & $\bullet$ \\
Angular & max\_norm\_L31/mode\_1/angle\_150 & $0.090$ & $6.26$ \scriptsize$[6.07,\,6.46]$ & $1.00$ & $0.00$ &  \\
Angular & max\_norm\_L31/mode\_1/angle\_270 & $0.237$ & $7.06$ \scriptsize$[6.86,\,7.23]$ & $1.00$ & $0.01$ &  \\
Angular & max\_norm\_L31/mode\_0/angle\_270 & $0.243$ & $6.87$ \scriptsize$[6.65,\,7.06]$ & $1.00$ & $0.00$ &  \\
Angular & max\_sim\_L23/mode\_0/angle\_270 & $0.377$ & $7.08$ \scriptsize$[6.87,\,7.26]$ & $1.00$ & $0.00$ &  \\
Angular & max\_norm\_L31/mode\_0/angle\_120 & $0.384$ & $7.29$ \scriptsize$[7.12,\,7.45]$ & $1.00$ & $0.00$ &  \\
Angular & max\_sim\_L23/mode\_1/angle\_270 & $0.418$ & $7.42$ \scriptsize$[7.25,\,7.58]$ & $1.00$ & $0.00$ &  \\
Angular & max\_norm\_L31/mode\_1/angle\_120 & $0.526$ & $7.49$ \scriptsize$[7.35,\,7.62]$ & $1.00$ & $0.00$ &  \\
Angular & max\_sim\_L23/mode\_0/angle\_60 & $0.599$ & $7.70$ \scriptsize$[7.58,\,7.80]$ & $1.00$ & $0.00$ &  \\
Angular & max\_sim\_L23/mode\_0/angle\_0 & $0.614$ & $7.88$ \scriptsize$[7.76,\,7.98]$ & $1.00$ & $0.00$ &  \\
Angular & max\_norm\_L31/mode\_1/angle\_0 & $0.629$ & $7.95$ \scriptsize$[7.84,\,8.05]$ & $1.00$ & $0.00$ &  \\
Angular & max\_sim\_L23/mode\_1/angle\_60 & $0.633$ & $7.79$ \scriptsize$[7.66,\,7.91]$ & $1.00$ & $0.00$ &  \\
Angular & max\_norm\_L31/mode\_1/angle\_60 & $0.651$ & $7.97$ \scriptsize$[7.87,\,8.06]$ & $1.00$ & $0.00$ &  \\
Angular & max\_norm\_L31/mode\_0/angle\_0 & $0.670$ & $7.92$ \scriptsize$[7.82,\,8.00]$ & $1.00$ & $0.00$ &  \\
Angular & max\_norm\_L31/mode\_0/angle\_60 & $0.672$ & $7.87$ \scriptsize$[7.78,\,7.95]$ & $1.00$ & $0.00$ &  \\
Angular & max\_sim\_L23/mode\_1/angle\_0 & $0.703$ & $7.98$ \scriptsize$[7.87,\,8.07]$ & $1.00$ & $0.00$ &  \\
CAA & L14/alpha\_2.0 & $-0.321$ & $5.96$ \scriptsize$[5.74,\,6.17]$ & $1.00$ & $0.00$ & $\bullet$ \\
CAA & L14/alpha\_4.0 & $-0.249$ & $1.24$ \scriptsize$[1.06,\,1.40]$ & $1.00$ & $0.05$ & $\bullet$ \\
CAA & L14/alpha\_6.0 & $0.000$ & $3.02$ \scriptsize$[2.72,\,3.32]$ & $0.99$ & $0.60$ & $\bullet$ \\
CAA & L14/alpha\_8.0 & $0.000$ & $2.90$ \scriptsize$[2.62,\,3.20]$ & $0.84$ & $0.56$ & $\bullet$ \\
CAA & L14/alpha\_1.0 & $0.146$ & $7.29$ \scriptsize$[7.11,\,7.46]$ & $1.00$ & $0.00$ &  \\
CAA & L14/alpha\_0.5 & $0.433$ & $7.70$ \scriptsize$[7.58,\,7.82]$ & $1.00$ & $0.00$ &  \\
\bottomrule
\end{tabular}
\caption{All sweep points, part 1 of 3.}
\label{tab:allsweep1}
\end{table*}

\begin{table*}[t]
\centering\small
\begin{tabular}{llrrrrc}
\toprule
Method & Configuration & Med. & Quality \scriptsize[95\% CI] & Uniq & Fail & Deg. \\
\midrule
CAA & L14/alpha\_0.0 & $0.712$ & $7.87$ \scriptsize$[7.74,\,7.98]$ & $1.00$ & $0.00$ &  \\
CAD mean-diff & alpha15.0 & $-0.183$ & $2.52$ \scriptsize$[2.28,\,2.77]$ & $1.00$ & $0.28$ & $\bullet$ \\
CAD mean-diff & alpha10.0 & $-0.180$ & $6.60$ \scriptsize$[6.41,\,6.80]$ & $1.00$ & $0.00$ & $\bullet$ \\
CAD mean-diff & alpha9.0 & $-0.107$ & $6.96$ \scriptsize$[6.78,\,7.13]$ & $1.00$ & $0.00$ & $\bullet$ \\
CAD mean-diff & alpha20.0 & $-0.035$ & $1.81$ \scriptsize$[1.52,\,2.10]$ & $0.99$ & $0.30$ & $\bullet$ \\
CAD mean-diff & alpha8.0 & $0.077$ & $7.20$ \scriptsize$[7.03,\,7.38]$ & $1.00$ & $0.00$ &  \\
CAD mean-diff & alpha7.0 & $0.139$ & $7.47$ \scriptsize$[7.31,\,7.61]$ & $1.00$ & $0.00$ &  \\
CAD mean-diff & alpha5.0 & $0.338$ & $7.70$ \scriptsize$[7.56,\,7.83]$ & $1.00$ & $0.00$ &  \\
CAD mean-diff & alpha3.0 & $0.402$ & $7.82$ \scriptsize$[7.70,\,7.92]$ & $1.00$ & $0.00$ &  \\
CAD mean-diff & alpha2.0 & $0.466$ & $7.84$ \scriptsize$[7.72,\,7.95]$ & $1.00$ & $0.00$ &  \\
CAD mean-diff & alpha1.0 & $0.583$ & $7.92$ \scriptsize$[7.81,\,8.01]$ & $1.00$ & $0.00$ &  \\
CAD probe & alpha13.0 & $-0.186$ & $4.34$ \scriptsize$[4.11,\,4.57]$ & $1.00$ & $0.04$ & $\bullet$ \\
CAD probe & alpha11.0 & $-0.149$ & $5.64$ \scriptsize$[5.44,\,5.83]$ & $1.00$ & $0.02$ & $\bullet$ \\
CAD probe & alpha15.0 & $-0.120$ & $2.60$ \scriptsize$[2.37,\,2.83]$ & $0.99$ & $0.20$ & $\bullet$ \\
CAD probe & alpha17.0 & $-0.048$ & $2.27$ \scriptsize$[1.97,\,2.56]$ & $0.99$ & $0.38$ & $\bullet$ \\
CAD probe & alpha9.0 & $0.034$ & $7.02$ \scriptsize$[6.85,\,7.19]$ & $1.00$ & $0.00$ &  \\
CAD probe & alpha7.0 & $0.268$ & $7.44$ \scriptsize$[7.27,\,7.59]$ & $1.00$ & $0.00$ &  \\
CAD probe & alpha5.0 & $0.366$ & $7.60$ \scriptsize$[7.44,\,7.74]$ & $1.00$ & $0.00$ &  \\
CAD probe & alpha3.0 & $0.437$ & $7.87$ \scriptsize$[7.75,\,7.99]$ & $1.00$ & $0.00$ &  \\
CAD probe & alpha1.0 & $0.569$ & $7.94$ \scriptsize$[7.84,\,8.04]$ & $1.00$ & $0.00$ &  \\
CAD projection & alpha9.0 & $-0.114$ & $1.98$ \scriptsize$[1.66,\,2.29]$ & $0.94$ & $0.39$ & $\bullet$ \\
CAD projection & alpha7.0 & $-0.049$ & $1.95$ \scriptsize$[1.65,\,2.24]$ & $0.96$ & $0.38$ & $\bullet$ \\
CAD projection & alpha5.0 & $-0.011$ & $2.16$ \scriptsize$[1.86,\,2.43]$ & $0.80$ & $0.40$ & $\bullet$ \\
CAD projection & alpha3.0 & $-0.002$ & $3.54$ \scriptsize$[3.25,\,3.80]$ & $0.83$ & $0.70$ & $\bullet$ \\
CAD projection & alpha2.0 & $0.181$ & $6.98$ \scriptsize$[6.78,\,7.17]$ & $1.00$ & $0.00$ &  \\
CAD projection & alpha1.0 & $0.199$ & $7.59$ \scriptsize$[7.44,\,7.73]$ & $1.00$ & $0.01$ &  \\
CAD projection & alpha0.5 & $0.371$ & $7.79$ \scriptsize$[7.67,\,7.90]$ & $1.00$ & $0.01$ &  \\
CAD projection & alpha0.3 & $0.444$ & $7.85$ \scriptsize$[7.74,\,7.96]$ & $1.00$ & $0.00$ &  \\
CAD projection & alpha0.1 & $0.580$ & $7.97$ \scriptsize$[7.88,\,8.06]$ & $1.00$ & $0.00$ &  \\
FairSteer & L29/thresh\_0.7/alpha\_4.0 & $-0.018$ & $2.08$ \scriptsize$[1.80,\,2.38]$ & $0.97$ & $0.36$ & $\bullet$ \\
\bottomrule
\end{tabular}
\caption{All sweep points, part 2 of 3.}
\label{tab:allsweep2}
\end{table*}

\begin{table*}[t]
\centering\small
\begin{tabular}{llrrrrc}
\toprule
Method & Configuration & Med. & Quality \scriptsize[95\% CI] & Uniq & Fail & Deg. \\
\midrule
FairSteer & L29/thresh\_0.5/alpha\_4.0 & $-0.012$ & $2.31$ \scriptsize$[2.02,\,2.61]$ & $0.94$ & $0.38$ & $\bullet$ \\
FairSteer & L29/thresh\_0.3/alpha\_4.0 & $-0.002$ & $2.31$ \scriptsize$[2.02,\,2.62]$ & $0.95$ & $0.39$ & $\bullet$ \\
FairSteer & L29/thresh\_0.3/alpha\_6.0 & $0.000$ & $2.42$ \scriptsize$[2.12,\,2.73]$ & $0.30$ & $0.40$ & $\bullet$ \\
FairSteer & L29/thresh\_0.3/alpha\_8.0 & $0.000$ & $4.28$ \scriptsize$[4.08,\,4.48]$ & $0.30$ & $0.82$ & $\bullet$ \\
FairSteer & L29/thresh\_0.5/alpha\_6.0 & $0.000$ & $2.38$ \scriptsize$[2.08,\,2.72]$ & $0.27$ & $0.40$ & $\bullet$ \\
FairSteer & L29/thresh\_0.5/alpha\_8.0 & $0.000$ & $4.34$ \scriptsize$[4.15,\,4.53]$ & $0.26$ & $0.83$ & $\bullet$ \\
FairSteer & L29/thresh\_0.7/alpha\_6.0 & $0.000$ & $2.38$ \scriptsize$[2.09,\,2.70]$ & $0.26$ & $0.39$ & $\bullet$ \\
FairSteer & L29/thresh\_0.7/alpha\_8.0 & $0.000$ & $4.29$ \scriptsize$[4.09,\,4.48]$ & $0.25$ & $0.82$ & $\bullet$ \\
FairSteer & L29/thresh\_0.3/alpha\_2.0 & $0.144$ & $6.65$ \scriptsize$[6.46,\,6.83]$ & $1.00$ & $0.01$ &  \\
FairSteer & L29/thresh\_0.5/alpha\_2.0 & $0.184$ & $6.68$ \scriptsize$[6.50,\,6.87]$ & $1.00$ & $0.01$ &  \\
FairSteer & L29/thresh\_0.7/alpha\_2.0 & $0.207$ & $6.74$ \scriptsize$[6.55,\,6.92]$ & $1.00$ & $0.00$ &  \\
FairSteer & L29/thresh\_0.3/alpha\_1.0 & $0.476$ & $7.80$ \scriptsize$[7.68,\,7.91]$ & $1.00$ & $0.00$ &  \\
FairSteer & L29/thresh\_0.5/alpha\_1.0 & $0.477$ & $7.80$ \scriptsize$[7.69,\,7.91]$ & $1.00$ & $0.00$ &  \\
FairSteer & L29/thresh\_0.5/alpha\_0.5 & $0.539$ & $7.90$ \scriptsize$[7.78,\,8.00]$ & $1.00$ & $0.00$ &  \\
FairSteer & L29/thresh\_0.7/alpha\_1.0 & $0.566$ & $7.77$ \scriptsize$[7.64,\,7.89]$ & $1.00$ & $0.00$ &  \\
FairSteer & L29/thresh\_0.7/alpha\_0.5 & $0.568$ & $7.83$ \scriptsize$[7.70,\,7.94]$ & $1.00$ & $0.00$ &  \\
FairSteer & L29/thresh\_0.3/alpha\_0.5 & $0.589$ & $7.89$ \scriptsize$[7.79,\,8.00]$ & $1.00$ & $0.00$ &  \\
FairSteer & L29/thresh\_0.7/alpha\_0.0 & $0.650$ & $7.93$ \scriptsize$[7.83,\,8.02]$ & $1.00$ & $0.00$ &  \\
FairSteer & L29/thresh\_0.5/alpha\_0.0 & $0.663$ & $7.93$ \scriptsize$[7.82,\,8.03]$ & $1.00$ & $0.00$ &  \\
FairSteer & L29/thresh\_0.3/alpha\_0.0 & $0.712$ & $7.87$ \scriptsize$[7.74,\,7.98]$ & $1.00$ & $0.00$ &  \\
SADI & strength\_40.0 & $-0.011$ & $2.98$ \scriptsize$[2.72,\,3.25]$ & $0.99$ & $0.40$ & $\bullet$ \\
SADI & strength\_20.0 & $0.000$ & $2.69$ \scriptsize$[2.40,\,2.96]$ & $0.98$ & $0.37$ & $\bullet$ \\
SADI & strength\_80.0 & $0.000$ & $3.28$ \scriptsize$[3.02,\,3.53]$ & $1.00$ & $0.53$ & $\bullet$ \\
SADI & strength\_10.0 & $0.119$ & $3.19$ \scriptsize$[2.97,\,3.42]$ & $1.00$ & $0.36$ & $\bullet$ \\
SADI & strength\_5.0 & $0.593$ & $7.53$ \scriptsize$[7.37,\,7.68]$ & $1.00$ & $0.01$ &  \\
SADI & strength\_2.0 & $0.596$ & $7.96$ \scriptsize$[7.86,\,8.05]$ & $1.00$ & $0.00$ &  \\
SADI & strength\_1.0 & $0.656$ & $7.91$ \scriptsize$[7.80,\,8.02]$ & $1.00$ & $0.00$ &  \\
\bottomrule
\end{tabular}
\caption{All sweep points, part 3 of 3.}
\label{tab:allsweep3}
\end{table*}

